%!TEX root = ../../thesis.tex
\section{目的}
\label{sec:purpose}

本研究の目的は，目標方向の決定をシナリオが担い，各区間の走行を基盤モデルが担うシステムにより，走行対象の環境での事前走行や再学習，自己位置推定を必要とせずに，シナリオに従って目的地まで移動できるかを明らかにすることである．
そのために，シミュレータ上の単一の棟で収集したデータのみでファインチューニングしたシステムを，データを収集していない別の棟で走行させて検証する．

本研究では，次の2点に取り組む．
\begin{itemize}
  \item シナリオ駆動型ナビゲーションの通路分類と経路追従に基盤モデルを応用し，事前学習で獲得された表現を保つように学習するシステムを構築する．
  \item 学習環境と，データを収集していない未学習環境の両方でシナリオに従った走行を行い，目的地に到達できるかを評価する．
\end{itemize}

% TODO: 修論での目的・貢献を書き足す（実環境での評価を含めるかなど）．
